\chapter[2021 --- Julius et al.]{2021: David Julius, Ardem Patapoutian}
\label{ch:2021_Julius}

\section*{Main Contribution}

\textit{Discovered receptors for temperature and touch (TRPV1, Piezo channels), explaining how the body senses physical stimuli---opening new avenues for pain therapy.}

\section{Official Citation}

for their discoveries of receptors for temperature and touch

\section{Background and Historical Context}

The 2021 Nobel Prize in Physiology or Medicine recognized David Julius and Ardem Patapoutian ``for their discoveries of receptors for temperature and touch.'' Their work identified the molecular sensors that allow us to perceive the physical world---the ion channels that detect painful heat, innocuous warmth, and mechanical force---and opened up new avenues for the development of novel pain therapies and a deeper understanding of how the nervous system processes sensory information.

The ability to perceive temperature and touch is fundamental to the survival of virtually all animals. Temperature perception enables organisms to detect environmental hazards, such as extreme heat or cold, and to regulate their body temperature. Touch perception enables organisms to navigate their environment, detect physical contact, and respond to mechanical stimuli such as pressure, vibration, and stretch. The molecular mechanisms by which sensory neurons detect these physical stimuli and convert them into electrical signals---a process known as sensory transduction---had been a central question in neuroscience for more than a century.

The concept of sensory receptors---specialized molecules that detect specific stimuli and initiate electrical signals in sensory neurons---dates back to the work of the German physiologist Max von Frey in the 1890s. Von Frey proposed that different types of sensory receptors were responsible for detecting different types of stimuli, including temperature, touch, pain, and pressure. However, the molecular identity of these receptors remained unknown for more than a century.

The development of molecular biology and genetics in the late twentieth century provided new tools for identifying sensory receptors. The cloning of the first ion channel---the \textit{Shaker} potassium channel in \textit{Drosophila}---in the 1980s demonstrated that ion channels could be identified and characterized using molecular genetic approaches. Subsequent cloning of ion channels involved in sensory transduction, including the capsaicin receptor and the mechanosensitive channels, provided direct evidence for the molecular basis of sensory perception.

The study of sensory transduction had been advanced by electrophysiological techniques that could record the activity of individual sensory neurons in response to specific stimuli. These studies had identified different classes of sensory neurons with distinct response properties---thermoreceptors that responded to temperature changes, mechanoreceptors that responded to mechanical stimuli, and nociceptors that responded to noxious (painful) stimuli. However, the molecular identity of the receptors responsible for detecting these stimuli remained unknown. The challenge was to bridge the gap between the physiological characterization of sensory neurons and the molecular identification of the receptors that gave them their specific response properties.

\section{The Discovery/Contribution}

\textbf{David Julius:} Julius, an American molecular biologist born in 1955 in Brooklyn, New York, working at the University of California, San Francisco, made the foundational discovery of the receptor for capsaicin---the active component of chili peppers that produces a sensation of burning pain. Capsaicin had been used for centuries as a topical analgesic, and it was known to activate a specific subset of sensory neurons called nociceptors (pain-sensing neurons). However, the molecular receptor for capsaicin had not been identified.

Julius and his colleagues used an ingenious approach to identify the capsaicin receptor. They constructed a cDNA library from sensory neurons that were known to be sensitive to capsaicin, and they screened this library in cultured cells for the ability to confer capsaicin sensitivity. After screening approximately 16,000 clones, they identified a single cDNA that conferred capsaicin sensitivity to previously unresponsive cells. This cDNA encoded a novel ion channel, which they named TRPV1 (transient receptor potential vanilloid 1).

TRPV1 was a member of the transient receptor potential (TRP) family of ion channels, a large and diverse family of channels that play important roles in sensory transduction. The TRP family was named after the \textit{trp} channel first discovered in \textit{Drosophila} photoreceptors, and its members share a common structure with six transmembrane domains and a pore loop between the fifth and sixth domains. Julius showed that TRPV1 was a nonselective cation channel that, when activated, allowed the influx of calcium and sodium ions, depolarizing the cell and generating an electrical signal.

Julius showed that TRPV1 was activated not only by capsaicin but also by temperatures above 43\,$^{\circ}$C---the threshold temperature for painful heat---as well as by low pH (acidic conditions), which are associated with tissue damage and inflammation. This dual activation by capsaicin and heat explained why capsaicin produced a burning sensation: it activated the same receptor that normally detected painful heat. The convergence of multiple stimuli on a single receptor explained the phenomenon of sensory cross-modal activation and provided a molecular explanation for why capsaicin, acid, and heat all produce burning pain.

The discovery of TRPV1 was a breakthrough in our understanding of temperature perception. It demonstrated that a single molecular receptor could detect both a chemical stimulus (capsaicin) and a physical stimulus (heat), and it provided a molecular explanation for the phenomenon of chemesthesis---the ability of chemical compounds to produce sensations of temperature or pain. Julius and his colleagues subsequently identified additional members of the TRP family that detected other temperature ranges, including TRPV2 (activated by temperatures above 52\,$^{\circ}$C), TRPV3 and TRPV4 (activated by warm temperatures of 27-35\,$^{\circ}$C), and TRPM8 (activated by cool temperatures below 26\,$^{\circ}$C and menthol). Together, these TRP channels provided a molecular ``thermometer'' system that covered the full range of temperatures relevant to biological sensation.

The identification of TRPV1 as the capsaicin receptor also had immediate therapeutic implications. Capsaicin-based topical creams had been used for centuries as folk remedies for pain relief, and the identification of TRPV1 explained how they worked: by activating and then desensitizing TRPV1-expressing nociceptors. This understanding led to the development of high-concentration capsaicin patches (8\%) for the treatment of neuropathic pain, which work by initially activating TRPV1 and then causing a prolonged period of desensitization that provides sustained pain relief.

\textbf{Ardem Patapoutian:} Patapoutian, an Armenian-American molecular biologist born in 1967 in Los Angeles, California, working at the Scripps Research Institute in La Jolla, California, made the foundational discovery of the molecular receptors for touch and mechanical sensation. While the existence of mechanosensitive ion channels---channels that open in response to mechanical force---had been postulated for decades, their molecular identity had remained elusive.

Patapoutian and his colleagues used a similar screen-based approach to identify mechanosensitive channels. They constructed a cDNA library from sensory neurons that were known to be mechanosensitive, and they screened this library in cultured cells for the ability to confer mechanosensitivity. The screening method involved using a laser-triggered glass pipette to apply mechanical force to individual cells while simultaneously recording electrical activity using the patch-clamp technique. After screening approximately 72 clones, they identified a family of ion channels they named Piezo (from the Greek word for ``pressure''): Piezo1 and Piezo2.

Piezo2 was shown to be the principal mechanosensor in sensory neurons, responsible for detecting light touch, proprioception (the sense of body position), and baroreception (the sense of blood pressure). Patapoutian demonstrated that Piezo2 was expressed in sensory neurons of the dorsal root ganglia and trigeminal ganglia, particularly in the mechanoreceptors that innervate the skin. Mice lacking Piezo2 had significant deficits in light touch sensation---they could not feel light touches on their skin---and proprioception, as evidenced by their inability to sense the position of their limbs in space. These mice were unable to walk normally on a rotating rod and had difficulty coordinating their movements, demonstrating the essential role of Piezo2 in touch and proprioception. Remarkably, Piezo2-deficient mice also had defects in breathing, suggesting that Piezo2-mediated mechanosensation plays a role in monitoring lung inflation and regulating respiratory rhythm.

Piezo1 was shown to be a mechanosensitive channel expressed in non-neuronal cells, including red blood cells, endothelial cells, and osteoblasts. Piezo1 plays important roles in vascular development, red blood cell volume regulation, and bone formation. It functions as a mechanosensor in endothelial cells, detecting the shear stress exerted by blood flow on the vessel wall and transducing this mechanical signal into biochemical responses that regulate vascular tone and remodeling. Mutations in Piezo1 have been linked to hereditary xerocytosis, a condition characterized by dehydrated red blood cells, demonstrating the clinical importance of mechanosensitive ion channels.

The structural determination of Piezo1 and Piezo2 using cryo-electron microscopy, published in 2018, revealed a remarkably unique structure---a large, trimeric complex with a propeller-like shape, featuring three curved ``blade'' domains that surround a central pore. This structure, which is unlike any previously known ion channel, provided insights into the mechanism by which Piezo channels detect and respond to mechanical force. The curved blades are thought to act as force sensors, flattening in response to membrane tension and thereby opening the central pore to allow ion flow.

The discoveries of Julius and Patapoutian revealed that sensory transduction---the conversion of physical and chemical stimuli into electrical signals---is mediated by a small number of highly specialized ion channels. TRP channels detect temperature and chemical stimuli, while Piezo channels detect mechanical stimuli. These discoveries provided a molecular framework for understanding how the nervous system perceives the physical world and opened up new avenues for the development of novel pain therapies.

\section{Impact on Modern Medicine}

The discoveries of Julius and Patapoutian have had a significant impact on our understanding of sensory physiology and have opened up new therapeutic approaches for pain management and other conditions.

\textbf{Pain management:} The identification of TRPV1 and other TRP channels as sensors of painful stimuli has provided new targets for the development of analgesic drugs. TRPV1 antagonists have been developed and tested in clinical trials for the treatment of chronic pain conditions, including neuropathic pain, inflammatory pain, and cancer pain. While some early TRPV1 antagonists produced side effects related to impaired heat sensation, more selective agents are being developed that could provide pain relief without compromising normal temperature perception.

The understanding of TRP channel function has also informed the development of topical analgesics based on capsaicin. High-concentration capsaicin patches (8\%) have been approved for the treatment of neuropathic pain associated with postherpetic neuralgia, diabetic peripheral neuropathy, and other conditions. These patches work by desensitizing TRPV1-expressing nociceptors, providing sustained pain relief.

\textbf{Touch and proprioception disorders:} The identification of Piezo2 as the principal mechanosensor has provided insights into the pathogenesis of touch and proprioception disorders. Mutations in Piezo2 have been linked to distal arthrogryposis type 5, a condition characterized by joint contractures and impaired proprioception. The understanding of Piezo2 function may lead to new therapeutic strategies for these conditions and for other disorders of touch and proprioception.

\textbf{Cardiovascular medicine:} Piezo1 plays important roles in vascular biology, including the regulation of blood pressure and vascular development. The identification of Piezo1 as a mechanosensor in endothelial cells has provided insights into the mechanisms of shear stress sensing and vascular remodeling. Piezo1 modulators are being explored as potential treatments for hypertension and other cardiovascular diseases.

\textbf{Bone biology:} Piezo1 is expressed in osteoblasts and plays a role in bone formation in response to mechanical loading. The understanding of Piezo1 function in bone cells has provided insights into the mechanisms of mechanotransduction in bone and has suggested potential therapeutic strategies for osteoporosis and other bone disorders.

\textbf{Sensory biology and neuroscience:} The discoveries of Julius and Patapoutian have transformed our understanding of sensory biology and have provided a molecular framework for studying how the nervous system processes sensory information. The identification of the molecular sensors for temperature, touch, and pain has enabled new approaches to studying sensory circuits, sensory coding, and the neural basis of perception.

David Julius and Ardem Patapoutian were awarded the Nobel Prize in Physiology or Medicine in 2021 ``for their discoveries of receptors for temperature and touch.'' Their work has fundamentally changed our understanding of sensory perception and has opened up new therapeutic avenues for pain management and other conditions.

\section{Legacy}

The legacy of Julius and Patapoutian's work extends across multiple areas of biology and medicine. Their discoveries have revealed the molecular sensors that allow us to perceive the physical world and have provided a foundation for understanding the mechanisms of sensory transduction.

David Julius's discovery of TRPV1 and other TRP channels has illuminated the molecular basis of temperature perception and has revealed the remarkable versatility of TRP channels as sensors of diverse stimuli, including temperature, chemicals, and mechanical force. The TRP channel family has emerged as one of a major and diverse families of ion channels in the body, with roles in sensory perception, thermoregulation, pain signaling, and many other physiological processes.

Ardem Patapoutian's discovery of the Piezo channels has revealed the molecular basis of touch and proprioception and has opened up a new field of mechanobiology---the study of how cells detect and respond to mechanical forces. The Piezo channels have been shown to play important roles not only in sensory neurons but also in non-neuronal cells, including red blood cells, endothelial cells, and osteoblasts, where they regulate a wide range of physiological processes.

Together, their work has provided a molecular framework for understanding how the nervous system perceives the physical world and has opened up new avenues for the development of novel pain therapies and treatments for sensory disorders. The legacy of their discoveries will continue to shape the future of sensory biology and medicine for generations to come.


\section{Modern Developments}

Julius and Patapoutian's sensory receptor work has led to modern pain research. Understanding of TRPV1 has led to development of TRPV1 antagonists for pain, though hyperthermia side effects remain challenging. Understanding of Piezo channels has revealed their role in touch, proprioception, and blood pressure regulation. Understanding of sensory neuron biology has informed development of channel blockers for neuropathic pain. Understanding of mechanotransduction has led to research on conditions like fibromyalgia and chronic pain.
